\documentclass[10pt,a4paper]{amsart}
\usepackage{amsmath,amssymb,amsthm,graphicx}
% Retained preamble for the Stacks Project chapter morphisms.tex (Morphisms of Schemes).
% The chapter file pinned-original-source.tex carries only \input{preamble} on its first line, so no
% preamble text lives inside the chapter itself.  The Stacks Project uses one single shared preamble for
% every chapter; the snapshot stages it verbatim in sources/stacks-chow/preamble.tex.  The block below is
% copied mechanically, line for line, from that file's lines 14-18 (the Xy-pic setup) and lines 154-217
% (the native theorem declarations together with the mathematical macros).  Nothing was added, removed or
% redefined.  The wrapper class and the amsmath/amssymb/amsthm/graphicx base packages are supplied by the
% assembler, so the class/\externaldocument part of the shared preamble is not repeated here.
% For commutative diagrams we use Xy-pic
\usepackage[all]{xy}

% We use 2cell for 2-commutative diagrams.
\xyoption{2cell}
% Theorem environments.
%
\theoremstyle{plain}
\newtheorem{theorem}[subsection]{Theorem}
\newtheorem{proposition}[subsection]{Proposition}
\newtheorem{lemma}[subsection]{Lemma}

\theoremstyle{definition}
\newtheorem{definition}[subsection]{Definition}
\newtheorem{example}[subsection]{Example}
\newtheorem{exercise}[subsection]{Exercise}
\newtheorem{situation}[subsection]{Situation}

\theoremstyle{remark}
\newtheorem{remark}[subsection]{Remark}
\newtheorem{remarks}[subsection]{Remarks}

\numberwithin{equation}{subsection}

% Macros
%
\def\lim{\mathop{\mathrm{lim}}\nolimits}
\def\colim{\mathop{\mathrm{colim}}\nolimits}
\def\Spec{\mathop{\mathrm{Spec}}}
\def\Hom{\mathop{\mathrm{Hom}}\nolimits}
\def\Ext{\mathop{\mathrm{Ext}}\nolimits}
\def\SheafHom{\mathop{\mathcal{H}\!\mathit{om}}\nolimits}
\def\SheafExt{\mathop{\mathcal{E}\!\mathit{xt}}\nolimits}
\def\Sch{\mathit{Sch}}
\def\Mor{\mathop{\mathrm{Mor}}\nolimits}
\def\Ob{\mathop{\mathrm{Ob}}\nolimits}
\def\Sh{\mathop{\mathit{Sh}}\nolimits}
\def\NL{\mathop{N\!L}\nolimits}
\def\CH{\mathop{\mathrm{CH}}\nolimits}
\def\proetale{{pro\text{-}\acute{e}tale}}
\def\etale{{\acute{e}tale}}
\def\QCoh{\mathit{QCoh}}
\def\Ker{\mathop{\mathrm{Ker}}}
\def\Im{\mathop{\mathrm{Im}}}
\def\Coker{\mathop{\mathrm{Coker}}}
\def\Coim{\mathop{\mathrm{Coim}}}

% Boxtimes
%
\DeclareMathSymbol{\boxtimes}{\mathbin}{AMSa}{"02}

%
% Macros for moduli stacks/spaces
%
\def\QCohstack{\mathcal{QC}\!\mathit{oh}}
\def\Cohstack{\mathcal{C}\!\mathit{oh}}
\def\Spacesstack{\mathcal{S}\!\mathit{paces}}
\def\Quotfunctor{\mathrm{Quot}}
\def\Hilbfunctor{\mathrm{Hilb}}
\def\Curvesstack{\mathcal{C}\!\mathit{urves}}
\def\Polarizedstack{\mathcal{P}\!\mathit{olarized}}
\def\Complexesstack{\mathcal{C}\!\mathit{omplexes}}
% \Pic is the operator that assigns to X its picard group, usage \Pic(X)
% \Picardstack_{X/B} denotes the Picard stack of X over B
% \Picardfunctor_{X/B} denotes the Picard functor of X over B
\def\Pic{\mathop{\mathrm{Pic}}\nolimits}
\def\Picardstack{\mathcal{P}\!\mathit{ic}}
\def\Picardfunctor{\mathrm{Pic}}
\def\Deformationcategory{\mathcal{D}\!\mathit{ef}}

% --- appended 2026-10-12 (Chief): the project own definitions of the reference/slogan/history
% display environments, copied mechanically from the shared Stacks preamble lines 8-12
% (\usepackage{verbatim} plus the three \newenvironment lines, which the chapter preamble keeps
% verbatim from the project preamble). Nothing added, removed or redefined: these environments
% carry no mathematical content and the project itself typesets them as comments. Without these
% three lines a counted statement whose own text uses them fails with Environment slogan undefined.
\usepackage{verbatim}
\newenvironment{reference}{\comment}{\endcomment}
\newenvironment{slogan}{\comment}{\endcomment}
\newenvironment{history}{\comment}{\endcomment}
\usepackage{cleveref}
\crefname{originallemma}{original lemma}{original lemmas}
\Crefname{originallemma}{Original lemma}{Original lemmas}
\makeatletter
\newlabel{lemma-locally-finite-type-characterize}{{01T2}{source}{}{sourcepointer-1}{}}
\newlabel{lemma-locally-finite-type-characterize@cref}{{[originallemma][0][]01T2}{[1][0][]source}{}{}{}}
\makeatother
\crefname{originallemma}{original lemma}{original lemmas}
\Crefname{originallemma}{Original lemma}{Original lemmas}
\makeatletter
\newlabel{properties-lemma-quasi-affine}{{01P9}{source}{}{sourcepointer-2}{}}
\newlabel{properties-lemma-quasi-affine@cref}{{[originallemma][0][]01P9}{[1][0][]source}{}{}{}}
\makeatother
\title{A quasi-affine scheme locally of finite type over a base admits an immersion into an affine space over that base}
\author{The Stacks Project Authors}
\date{stacks-project:morphisms.tex:local-snapshot-20261012; GNU FDL 1.2 or later}
\begin{document}
\maketitle
\noindent\textit{Editorial note.} All delivered mathematical text is the author's own, copied from the pinned Stacks Project chapter file pinned-original-source.tex; no statement or proof step was rewritten, reordered or removed.\par\medskip
\noindent\textit{Editorial note.} Following the multi-claim convention for this source, source\_url is the source-level URL https://stacks.math.columbia.edu/ and the source label lemma-quasi-affine-finite-type-over-S is carried in claim\_key and here rather than in a per-claim tag URL.\par\medskip
\noindent\textit{Editorial note.} The preamble is the chapter's shared Stacks Project preamble, retained as sources/stacks-morphisms/retained-preamble.tex: a line-for-line excerpt of the shared Stacks preamble staged with this snapshot. No macro was invented or redefined, and this card adds no preamble text of its own.\par\medskip
\noindent\textit{Editorial note.} Statement span 9372-9378 and proof span 9380-9410 come from the halfA pool card operations/pool-cards/stacks-morphisms-halfA.json and were re-read in the pinned chapter file before this card was written.\par\medskip
\noindent\textit{Editorial note.} This card was previously skipped because the shared gate substring-matched my chapter name against the published items of more-morphisms.tex; the Lead repaired the gate to exact file-name equality and the candidate was re-gated (RESULT=OK) before being opened.\par\medskip
\noindent\textit{Editorial note.} The proof names the coordinates of the affine space as the chart functions together with the chosen generators; the author's notation is preserved.\par\medskip
\noindent\textit{Editorial note.} 2 cross-references occur inside the counted spans; they are listed in external\_references as navigation pointers only.\par\medskip
\noindent\textit{Allowed standard prerequisites (editorial).}\begin{itemize}
\item A quasi-affine scheme is quasi-compact and identifies with an open subscheme of the spectrum of its ring of global functions, and the open sets where a global function is invertible and the corresponding localized ring is the ring of sections form a basis when they are affine. Cited at proof lines 9381-9386. Supplies the finite cover by standard affine opens used to build the immersion.
\item The same-chapter characterization of locally finite type morphisms by finite type ring maps on affine charts. Cited at proof lines 9389-9391. Makes each coordinate ring of a chart a finite type algebra over the coordinate ring of an affine open of the base.
\item The construction of a morphism into an affine space over the base from global sections of the structure sheaf and the identification of preimages of standard opens. Used at proof lines 9394-9405. Produces the candidate morphism and identifies the preimage of the standard open cut out by the j-th coordinate with the j-th chart.
\end{itemize}\medskip
\ifdefined\hypertarget\hypertarget{sourcepointer-1}{}\fi
\noindent\textit{Source reference.} Original lemma 01T2: focused-seeds/AGP/sources/stacks-snapshot-20261011/tags.txt, source lines 2359--2359. Same-chapter reference used in the retained author text at source line(s) 2717, 2739, 2801, 2863, 3019, 3035, 3157, 3348, 3613, 3707, 3868, 5197, 5282, 6702, 9325, 9390, 9440, 11421, 13422, 13686. It is supplied only as a navigation pointer so that the retained text compiles; the counted item delivers the author's own statement and the author's complete proof, and no part of the referenced item is copied in or substituted for the counted argument.\par\medskip
\ifdefined\hypertarget\hypertarget{sourcepointer-2}{}\fi
\noindent\textit{Source reference.} Original lemma 01P9: focused-seeds/AGP/sources/stacks-snapshot-20261011/tags.txt, source lines 2222--2222. Same-chapter reference used in the retained author text at source line(s) 2266, 9383, 9386. It is supplied only as a navigation pointer so that the retained text compiles; the counted item delivers the author's own statement and the author's complete proof, and no part of the referenced item is copied in or substituted for the counted argument.\par\medskip
\begin{lemma}
\label{lemma-quasi-affine-finite-type-over-S}
Let $\pi : X \to S$ be a morphism of schemes.
Assume that $X$ is quasi-affine and that $\pi$ is locally of finite type.
Then there exist $n \geq 0$ and an immersion $i : X \to \mathbf{A}^n_S$
over $S$.
\end{lemma}

\begin{proof}
Let $A = \Gamma(X, \mathcal{O}_X)$. By assumption $X$ is quasi-compact
and is identified with an open subscheme of $\Spec(A)$, see
Properties, Lemma \ref{properties-lemma-quasi-affine}.
Moreover, the set of opens $X_f$, for those $f \in A$ such that $X_f$ is
affine, forms a basis for the topology of $X$, see the proof of
Properties, Lemma \ref{properties-lemma-quasi-affine}.
Hence we can find a finite number of $f_j \in A$, $j = 1, \ldots, m$ such that
$X = \bigcup X_{f_j}$, and such that $\pi(X_{f_j}) \subset V_j$ for
some affine open $V_j \subset S$. By
Lemma \ref{lemma-locally-finite-type-characterize}
the ring maps $\mathcal{O}(V_j) \to \mathcal{O}(X_{f_j}) = A_{f_j}$
are of finite type. Thus we may choose $a_1, \ldots, a_N \in A$ such that
the elements $a_1, \ldots, a_N, 1/f_j$ generate
$A_{f_j}$ over $\mathcal{O}(V_j)$ for each $j$. Take $n = m + N$ and
let
$$
i : X \longrightarrow \mathbf{A}^n_S
$$
be the morphism given by the global sections
$f_1, \ldots, f_m, a_1, \ldots, a_N$ of the structure sheaf of $X$.
Let $D(x_j) \subset \mathbf{A}^n_S$ be the open subscheme where the
$j$th coordinate function is nonzero.
Then for $1 \leq j \leq m$ we have $i^{-1}(D(x_j)) = X_{f_j}$ and
the induced morphism $X_{f_j} \to D(x_j)$ factors through the affine
open $\Spec(\mathcal{O}(V_j)[x_1, \ldots, x_n, 1/x_j])$
of $D(x_j)$. Since the ring map
$\mathcal{O}(V_j)[x_1, \ldots, x_n, 1/x_j] \to A_{f_j}$ is
surjective by construction we conclude that $i^{-1}(D(x_j)) \to D(x_j)$
is an immersion as desired.
\end{proof}

\end{document}
